\chapter{从自旋筛选到概率振幅}
\label{ch:spin}

\section{Stern--Gerlach实验}
绪论中的光谱和碰撞实验从能量差推断原子有离散能级。若要进一步预言一次测量会落在哪个出口，
可以从一个只给两种读数的实验开始。银原子穿过磁体后分成两束，挡住一束就能用另一束制备下一次测量的样品。
把第二台磁体转过一个角度，我们既能数出每个出口有多少原子，也能问前一次筛选的信息还保留了多少。
这正是选自旋作起点的原因。只要改变保留哪束原子、下一台磁体朝向哪里，就能逐项检验制备和测量怎样决定结果的概率。

\subsection{装置与束流分裂}
\label{ss:sg-apparatus}

Stern--Gerlach实验要检验原子磁矩在一个给定方向上的投影是否连续。
加热银产生原子束，以狭缝准直，再让原子经过具有强梯度的磁场，
最后记录原子落点。选择中性原子使束流偏转主要来自磁矩与磁场梯度的耦合，
而不被净电荷的Lorentz力支配\cite{sg}。

\begin{figure}[htbp]
    \centering
    \tikzfigure[width=.74\linewidth]{sg-beam}
    \caption{银原子束从左向右传播。$\ell$ 标出梯度磁场区的长度，$D$ 是出磁体后的漂移距离；在右侧屏上沿 $z$ 方向量取两斑间距 $2\abs{z_\mathrm{det}}$。}
    \label{fig:sg-beam}
\end{figure}

均匀磁场可以改变磁矩的取向，却不给定磁矩一个持续的平移分离力。
空间分束需要场梯度。大学物理中的磁偶极能是 $-\bm\mu\cdot\bm B$。
量子理论保留这个空间向量的点积，把三个磁矩分量写成作用在原子内部态上的算符。
在束轴附近取局域磁场沿 $z$ 方向，即 $\bm B(z)\simeq B_z(z)\bm e_z$。
以 $\ket{s}$ 表示归一化的磁矩投影本征态，$\op\mu_z\ket{s}=\mu_z\ket{s}$，相应的相互作用能为
\begin{equation}
    \op H_{\interaction}(z)=-\hatbm{\mu}\cdot\bm B(z),\qquad
    E_{\interaction,s}(z)=\mel{s}{\op H_{\interaction}(z)}{s}
    \simeq-\mu_z B_z(z).
\end{equation}
点号把三个空间分量收缩成一个能量算符。外层尖括号是态空间里的内积，取的是 $\ket{s}$ 与 $\op H_{\interaction}(z)\ket{s}$ 的内积，
给出一个数。这个分支有确定的磁矩投影，取期望便得到该分支的能量。若原子保持在这一分支上，
就可以把它当作质心的有效势能。沿横向偏转坐标 $z$，有
\begin{equation}
    F_z=-\frac{\p E_{\interaction,s}}{\p z}
        \simeq\mu_z\frac{\p B_z}{\p z}.
\end{equation}
局域磁场方向给出测量磁矩分量的轴，称为量子化轴。该近似要求束流所在区域的场方向明确，
且原子穿过磁体时从一个内部能量分支跳到另一分支的概率很小。

为了把力与图\ref{fig:sg-beam}中屏上的位置联系起来，设磁体长度为$\ell$，
其后到探测面的自由飞行距离为$D$，原子质量为$M_\mathrm{Ag}$，
纵向速度为$v$，在窄束范围把梯度近似为常量$G=\p_zB_z$。
原子在磁体中停留$\tau=\ell/v$，离开磁体时的横向速度和位移分别为
\begin{subequations}
    \begin{alignat}{2}
        \Delta v_z&=\frac{\mu_zG}{M_\mathrm{Ag}}\tau,
        &\quad z_\mathrm{mag}&=\frac{\mu_zG}{2M_\mathrm{Ag}}\tau^2,\\
        z_\mathrm{drift}&=\Delta v_z\frac{D}{v},
        &\quad z_\mathrm{det}&=\frac{\mu_zG\ell}{M_\mathrm{Ag}v^2}
            \bigp{D+\ell/2}.
    \end{alignat}
\end{subequations}
因此在速度和几何条件固定时，探测位置直接标记$\mu_z$。
速度分布与狭缝宽度使每个分支具有有限宽度；
真正要比较的是连续的偏转带与分离的束流分支。

若磁矩只是大小固定而方向连续随机的经典小磁针，
$\mu_z=\mu\cos\theta$会在一个区间内连续取值。
银原子束却出现两个主要分支，这要求内部磁矩的投影有离散结构。
现代解释中，基态银原子的闭壳层不贡献净电子角动量，
剩余未配对电子的轨道角动量为零，主要磁矩来自自旋(spin)。
自旋是电子具有的内禀角动量，即使轨道角动量为零，磁场沿任何选定轴仍可分辨它的两个投影。
忽略此处不能分辨的核自旋结构，电子自旋投影为$\pm\hbar/2$，
磁矩约为$\hatbm{\mu}=-g_s\muB\hatbm{S}/\hbar$，其中 $g_s$ 为电子的无量纲自旋磁矩因子，约为 2。
磁矩与电子自旋方向相反，因此屏上“向上偏转”与“正自旋标签”的对应
还取决于梯度和电荷符号；实验标签应由所测算符统一规定。

\subsection{连续筛选与条件概率}
\label{ss:sg-sequence}

分束说明可能的读数，连续筛选则研究测量后保留下来的原子具有怎样的状态。
先沿一个方向筛选，再沿另一个方向筛选，最后回到原方向测量，能否重现第一次的确定读数？

把能由理想$z$分析器确定区分的两种内部态记作$\ket{+z}$和$\ket{-z}$，
测量自旋分量 $\op{S}_z$ 时，相应读数分别为 $+\hbar/2$ 和 $-\hbar/2$。筛选出$+z$后再做同方向测量，理想情况下仍得$+z$。
改用$x$方向测量会得到两种输出；若再筛出$+x$送入$z$分析器，$-z$输出重新出现。
改变分析器方向会改变测后态；预测后续计数时，需要同时描述测量概率和筛选后的制备。

\begin{figure}[htbp]
    \centering
    \tikzfigure{sg-filters}
    \caption{从左向右依次保留 $+z$ 束和 $+x$ 束。末端两支的计数之和就是进入最后一台分析器的原子数。}
\end{figure}

设进入第一次$z$筛选的原子数为$N$，
第一次保留$N_{+z}$个，第二次$x$筛选再保留$N_{+x,+z}$个。
最后$z$分析的两个计数记为$N_{+z,+x,+z}$和$N_{-z,+x,+z}$。
例如最后一步的条件频率为
\begin{equation}
    p(-z\mid+x,+z)
    \simeq\frac{N_{-z,+x,+z}}{N_{+x,+z}}.
\end{equation}
分母是实际进入最后分析器的原子数。
理想制备中，只要第二次筛选确实制备了同一个$+x$态，
最后的概率就由这个态决定，不再依赖更早的筛选记录。
下面的态与投影描述将把这一实验操作写成可计算的规则。

\section{态与测量概率}

筛出 $+z$ 束后，沿 $z$ 分析只有一种结果，沿 $x$ 分析却有两种。
若想把同一制备用到不同朝向的分析器上，就需要一种能对各个出口都给出概率的描述。
量子力学用\emph{态}描述制备，再把选定的分析方向代入，计算对应的读数分布。

\begin{definition}{纯态、振幅与整体相位}{qm-pure-state}
    纯态(pure state)用归一化(normalized)右矢$\ket{\psi}$表示，满足$\brkt{\psi}{\psi}=1$。
    $\ket{\psi}$与$\e^{\i\chi}\ket{\psi}$（$\chi$为实数）代表同一纯态。
    对某个归一化分析态$\ket{a}$，复数
    \begin{equation}
        c_a=\brkt{a}{\psi}
    \end{equation}
    称为概率振幅。
    $\bra{a}$作用于$\ket{\psi}$得到复数，而
    外积(outer product)$\ktbr{a}{b}$是算符，其作用为
    \begin{equation}
        (\ktbr{a}{b})\ket{\psi}=\ket{a}\brkt{b}{\psi}.
    \end{equation}
\end{definition}

把两个互相排斥、合起来包含全部输出的态用作一组基。沿用直角坐标中“沿一根轴取分量”的做法，
$\bra{+z}$ 取出 $+z$ 分量，对 $\ket{+z}$ 给 1，对 $\ket{-z}$ 给 0。两个基矢满足的这种关系称为正交归一(orthonormal)。
用 $s,s'\in\{+,-\}$ 表示出口的正负号，有
\begin{equation}
    \brkt{sz}{s'z}=\delta_{ss'},\qquad
    \one=\ktbr{+z}{+z}+\ktbr{-z}{-z}.
\end{equation}
一般纯态写成
\begin{equation}
    \ket{\psi}=c_+\ket{+z}+c_-\ket{-z},\qquad \abs{c_+}^2+\abs{c_-}^2=1.
\end{equation}
\begin{postulate}{Born规则}{qm-born}
    对归一化纯态 $\ket{\psi}$，在一组正交归一基中作理想测量，
    对应基矢 $\ket{a}$ 的出口概率为 $p(a\mid\psi)=\abs{\brkt{a}{\psi}}^2$。
\end{postulate}
例如 $p(\pm z)=\abs{c_\pm}^2$。实验可以在相同制备下转动分析器，逐一检验这个公设给出的分布。
归一化态乘同一个相位$\e^{\i\chi}$不会改变任何概率。两个系数都非零时，$c_-/c_+$的相位称为相对相位(relative phase)，可由实验读出。
左矢取共轭转置：$\bra{\psi}=c_+\conj\bra{+z}+c_-\conj\bra{-z}$。

例如$\brkt{\psi}{\psi}$展开为四项，其中两个交叉项因
$\brkt{+z}{-z}=0$而消失，留下$\abs{c_+}^2+\abs{c_-}^2$。
在有序基$\bigp{\ket{+z},\ket{-z}}$下，同一态及一个外积写成
\begin{subequations}
    \begin{alignat}{2}
        \ket{\psi}&\ \longleftrightarrow
        \begin{pmatrix}c_+\\c_-\end{pmatrix},
        &\quad \bra{\psi}&\ \longleftrightarrow
        \begin{pmatrix}c_+\conj&c_-\conj\end{pmatrix},\\
        \ktbr{+z}{-z}&\ \longleftrightarrow
        \begin{pmatrix}0&1\\0&0\end{pmatrix}.
    \end{alignat}
\end{subequations}
最后这个算符把$\ket{-z}$变成$\ket{+z}$，并把$\ket{+z}$映为零向量。
基矢本身对应可制备的态，系数列则是该态在所选基底下的表示。

怎样从计数确定 $x$ 分析器筛出的态？把 $+x$ 束送到 $z$ 分析器，两束计数各半。
由 Born 规则，两个 $z$ 分量的模都是 $1/\sqrt2$。可以给 $\ket{+x}$ 和基矢 $\ket{-z}$ 选择相位，
使这两个系数都是正实数。$-x$ 态要与 $+x$ 态正交，于是系数一正一负。
再将 $+y$ 束分别送入 $z$ 和 $x$ 分析器，两组计数也都各半。前一组把它的形式限制为
$(\ket{+z}+\e^{\i\varphi}\ket{-z})/\sqrt2$，后一组要求
$\abs{1+\e^{\i\varphi}}^2/4=1/2$，即 $\cos\varphi=0$。与右手坐标系一致地取 $+y$ 对应 $\varphi=\pi/2$，
便有
\begin{equation}
    \ket{\pm x}=\frac{\ket{+z}\pm\ket{-z}}{\sqrt2},\qquad
    \ket{\pm y}=\frac{\ket{+z}\pm\i\ket{-z}}{\sqrt2}.
\end{equation}
于是连续筛选的条件概率(conditional probability)为
$\abs{\brkt{+x}{+z}}^2=1/2$及$\abs{\brkt{\pm z}{+x}}^2=1/2$。
若分母是最初已经制备为$+z$的束流，最后每一路只占$1/4$，其余被$x$筛选器挡去。

\begin{example}{逐次筛选的状态与计数}{qm-filter-chain}
    从归一化的$\ket{+z}$出发，$+x$筛选对应算符
    $\op{P}_{+x}=\ktbr{+x}{+x}$。作用后的未归一化态是
    \begin{equation}
        \op{P}_{+x}\ket{+z}
        =\ket{+x}\brkt{+x}{+z}
        =\frac1{\sqrt2}\ket{+x}.
    \end{equation}
    范数平方$1/2$给出通过率；除以范数后，保留束的态为$\ket{+x}$。
    再投影到最后的$\pm z$输出，有
    \begin{equation}
        \op{P}_{\pm z}\op{P}_{+x}\ket{+z}=\frac12\ket{\pm z}.
    \end{equation}
    两个范数平方均为$1/4$，这是相对于最初$+z$束的联合概率。
    除以第二步通过率$1/2$，得到最后一步的两个条件概率均为$1/2$。
    筛选后的归一化与通过率分别保存，才能在多级装置中正确计算计数。
\end{example}

\section{Pauli算符与任意方向}
已经知道沿 $x,y,z$ 筛选出的态。若把磁体停在 $x$ 与 $z$ 之间，怎样从旋转的角度算出两个出口的通过率？
可以仿照经典角动量取分量的方法，把三个方向的读数写成一个矢量算符的三个分量，再沿磁体方向投影。

\begin{definition}{自旋二分之一的Pauli算符}{qm-pauli}
    把方向 $i$ 的两个自旋读数以 $\hbar/2$ 为单位，标为 $+1$ 和 $-1$。
    相应的无量纲 Pauli 算符为 $\op{\sigma}_i=\ktbr{+i}{+i}-\ktbr{-i}{-i}$，其中 $i=x,y,z$。
    自旋沿该轴的分量是 $\op{S}_i=(\hbar/2)\op{\sigma}_i$。
\end{definition}
例如将上节求出的两个 $x$ 态代入 $\op{\sigma}_x$，对角项相消，交叉项相加。
三个方向都改用 $z$ 基表示，得到
\begin{subequations}
    \begin{alignat}{1}
        \op{\sigma}_z&=\ktbr{+z}{+z}-\ktbr{-z}{-z},\\
        \op{\sigma}_x&=\ktbr{+z}{-z}+\ktbr{-z}{+z},\\
        \op{\sigma}_y&=-\i\ktbr{+z}{-z}+\i\ktbr{-z}{+z}.
    \end{alignat}
\end{subequations}
由外积相乘直接得到
\begin{equation}
    \op{\sigma}_i\op{\sigma}_j=\delta_{ij}\one+\i\epsilon_{ijk}\op{\sigma}_k.
\end{equation}

这里$\hatbm{\sigma}=(\op{\sigma}_x,\op{\sigma}_y,\op{\sigma}_z)$是三个Pauli算符组成的向量，
$\hatbm{S}$是自旋角动量，量纲为作用量。
例如由外积乘法，
\begin{alignat*}{1}
    \op{\sigma}_x\op{\sigma}_y &=\i\ktbr{+z}{+z}-\i\ktbr{-z}{-z}
        =\i\op{\sigma}_z,\\
    \op{\sigma}_y\op{\sigma}_x&=-\i\op{\sigma}_z.
\end{alignat*}
因此$[\op{S}_x,\op{S}_y]=\i\hbar\op{S}_z$。
同样相乘可得$\op{\sigma}_i^2=\one$，由此构成全部Pauli乘法规则。

把 $\op{\sigma}_z$ 作用在 $\ket{+z}$ 上，得到 $+\ket{+z}$，作用在 $\ket{-z}$ 上则得到 $-\ket{-z}$。
一个算符在某个非零态上只给出数乘，便称这个态为它的本征态，所乘的数为本征值；
在相应本征态上测量这个自旋分量，以概率一得到该本征值。
任意单位向量$\bm n$对应$\op{S}_{\bm n}=(\hbar/2)\bm n\cdot\hatbm{\sigma}$，
因为$(\bm n\cdot\hatbm{\sigma})^2=\one$，其本征值仍为$\pm\hbar/2$。

若$\bm n=(\sin\theta\cos\phi,\sin\theta\sin\phi,\cos\theta)$，则
\begin{equation}
    \ket{+\bm n}=\cos\frac\theta2\ket{+z}
    +\e^{\i\phi}\sin\frac\theta2\ket{-z}.
    \label{eq:spinor}
\end{equation}

这个式子可直接由本征方程求出。在$z$基中，
\begin{equation}
    \bm n\cdot\hatbm{\sigma}=
    \begin{pmatrix}
        \cos\theta&\e^{-\i\phi}\sin\theta\\
        \e^{\i\phi}\sin\theta&-\cos\theta
    \end{pmatrix}.
\end{equation}
令正本征态的系数为$(a,b)\transp$，第一行给出
$(1-\cos\theta)a=\e^{-\i\phi}\sin\theta\,b$。
在$0<\theta<\pi$时，
\begin{equation}
    \frac ba=\e^{\i\phi}\frac{1-\cos\theta}{\sin\theta}
    =\e^{\i\phi}\tan\frac\theta2.
\end{equation}
选$a$为实数并归一化，得到$a=\cos(\theta/2)$、
$b=\e^{\i\phi}\sin(\theta/2)$；两极取连续极限。
这里$\theta$是$\bm n$与$z$轴的夹角，$\phi$是其方位角。

投影到$+\bm n$的算符为
\begin{equation}
    \op{P}_{+\bm n}=\ktbr{+\bm n}{+\bm n}
    =\frac{\one+\bm n\cdot\hatbm{\sigma}}{2}.
\end{equation}
有限维算符 $\op{A}$ 的迹(trace)等于它在任一正交归一基中对角矩阵元之和，记为 $\tr\op{A}$。由于 $\tr(\ktbr{a}{b})=\brkt{b}{a}$，
两个一维投影乘积的迹正好给出 $\abs{\brkt{+\bm m}{+\bm n}}^2$。
由$\tr\op{\sigma}_i=0$和$\tr(\op{\sigma}_i\op{\sigma}_j)=2\delta_{ij}$，
$\tr(\op{P}_{+\bm m}\op{P}_{+\bm n})=(1+\bm m\cdot\bm n)/2$。
两个方向间的夹角为$\gamma$时，
$p(+\bm m\mid+\bm n)=(1+\bm m\cdot\bm n)/2=\cos^2(\gamma/2)$。
转动分析器改变了两个自旋态的重叠（图\ref{fig:spin-angle-readout}）。若将每个自旋纯态
$\ket{+\bm n}$ 画在单位球面方向 $\bm n$ 上，就得到 Bloch 球；沿同一轴的两种确定输出在球上相对。
光子线偏振筛选中，透射概率为 $\cos^2\vartheta$，其中 $\vartheta$ 是两根偏振轴的较小夹角（$0$ 至 $\pi/2$）。
与自旋形式比较，相应的 Bloch 球角距为 $2\vartheta$。

\begin{figure}[htbp]
    \centering
    \tikzfigure{spin-angle-readout}
    \caption{固定制备 $\ket{+z}$。左图从 $z$ 轴量取分析轴的转角 $\gamma$；右图实线、虚线分别给出正、负出口的计数比例，交点对应 $\gamma=\pi/2$。}
    \label{fig:spin-angle-readout}
\end{figure}

\section{单粒子干涉}
\label{sec:interference}
$+x$ 和 $-x$ 两种制备沿 $z$ 测量时都给出各半的计数，换到 $x$ 方向分析却走向相反出口。
造成这个差别的是两个 $z$ 振幅之间的相对相位。
把两个自旋标签换成粒子可走的两条路径，我们就能直接在一路上调节相位，
再让两路通向同一个探测器，观察计数怎样随相位变化。

电子衍射实验测量电子束被晶体散射后的角分布，
强度峰随电子动量和晶体取向改变；外村彰（Akira Tonomura）等人的单电子干涉实验则降低入射率，
逐个记录位置，随着事件累积形成稳定条纹\cite{davisson,tonomura}。
两类测量共同要求理论同时描述离散探测事件与依赖相位的统计分布。

在两路径实验中，先只开放路径$a$，记录每个输出的计数；
再只开放$b$重复测量；最后同时开放两路，并扫描一路上可控的相位。
若每次随机选择一路制备，输出概率就是两个单路概率的加权和；保持单路透过率时，这个和与相位无关。
实验条纹因此要求保留两路之间的相对相位信息。

用两个正交路径态$\ket{a},\ket{b}$表示刚经过分束器后的两种传播模式。
它们的正交性说明在该位置可以区分路径；合束器随后把这两个模式送入共同的输出模式，
从而使不同路径的振幅在同一探测结果上叠加。
制备相干叠加(coherent superposition)
$\ket{\psi}=(\ket{a}+\e^{\i\varphi}\ket{b})/\sqrt2$，再用能分析
$\ket{\pm}=(\ket{a}\pm\ket{b})/\sqrt2$的装置重合路径，得到
\begin{equation}
    p(+)=\frac{1+\cos\varphi}{2},\qquad p(-)=\frac{1-\cos\varphi}{2}.
    \label{eq:fringe}
\end{equation}

求这两个概率只需要内积：
\begin{subequations}
    \begin{alignat}{2}
        \brkt{+}{\psi} &=\frac{1+\e^{\i\varphi}}2,
        &\quad \abs{\brkt{+}{\psi}}^2 &=\frac{(1+\e^{\i\varphi})(1+\e^{-\i\varphi})}{4},\\
        \brkt{-}{\psi} &=\frac{1-\e^{\i\varphi}}2,
        &\quad \abs{\brkt{-}{\psi}}^2 &=\frac{(1-\e^{\i\varphi})(1-\e^{-\i\varphi})}{4}.
    \end{alignat}
\end{subequations}
两条正交路径各自的占据概率都是$1/2$，
相位却改变了合束后的输出。于是测量装置需要说明\emph{在哪个基底中读出}。

这是对两条路径的复振幅先相加再取模平方。
更一般地，在传播的中间时刻插入单位算符，并用中间态的一组正交归一基展开它，得到
\begin{equation}
    \mel{f}{\op{U}_2\op{U}_1}{i}
    =\sum_a\mel{f}{\op{U}_2}{a}\mel{a}{\op{U}_1}{i}.
\end{equation}
在这里对路径的振幅求和，描述粒子从制备端到探测端的连续传播。

若装置留下正交的路径记录，或者每次随机化相位且取平均，交叉项消失，
两输出概率均为$1/2$。逐个发射粒子时，只要两路保持相干，积累的计数仍会显示条纹。
\chapref{ch:entanglement}将用粒子与路径记录器的联合态计算条纹可见度。

\checkpoint{纯态
    $(\ket{+z}+\e^{\i\varphi}\ket{-z})/\sqrt2$
    的$z$测量概率与$\varphi$无关。怎样判断这个相位是否可测？}{
    先将测量基换为$x$方向。因为
    $\bra{+x}=(\bra{+z}+\bra{-z})/\sqrt2$，
    \begin{equation}
        \brkt{+x}{\psi}=\frac{1+\e^{\i\varphi}}2,
        \qquad
        p(+x)=\frac{1+\cos\varphi}{2}.
    \end{equation}
    再换为$y$方向。
    $\ket{+y}=(\ket{+z}+\i\ket{-z})/\sqrt2$的左矢必须共轭复系数，
    因而
    \begin{equation}
        \brkt{+y}{\psi}=\frac{1-\i\e^{\i\varphi}}2,
        \qquad
        p(+y)=\frac{1+\sin\varphi}{2}.
    \end{equation}
    两组独立计数分别读出$\cos\varphi$和$\sin\varphi$，
    可以确定$\varphi$模$2\pi$的值。
    单测$x$方向时，$\varphi$与$-\varphi$给出相同的读数分布。

    对一般两态纯态$c_+\ket{+z}+c_-\ket{-z}$，
    三个分析方向给出
    \begin{equation}
        \ave{\op{\sigma}_x}=2\Re(c_+\conj c_-),\quad
        \ave{\op{\sigma}_y}=2\Im(c_+\conj c_-),\quad
        \ave{\op{\sigma}_z}=\abs{c_+}^2-\abs{c_-}^2.
    \end{equation}
    三组计数给出两个分量的模，以及它们的干涉项 $c_+\conj c_-$。
    同时将$c_+,c_-$乘$\e^{\i\chi}$，这些量保持不变，因此测量概率与整体相位(global phase)无关。}

\reading{MIT 8.05讲义2核对自旋态和振幅规则\cite{mitspin}。
    接下来要把两输出推广到任意可观测量，并明确无穷维情况下“算符是什么”。}
